problem:  
Let \((M^{n+1}, g)\) be a smooth, globally hyperbolic Lorentzian manifold satisfying the **null energy condition**, and let \(\Sigma\subset M\) be a spacelike Cauchy hypersurface containing a compact, two‑sided **stable marginally outer trapped surface** \(S_0\subset\Sigma\) with \(\theta_+(S_0)=0\).  

Let \(\mathcal{N}^+\) be the null hypersurface generated by future‑directed affinely parameterized null geodesics \(\ell_+\) orthogonal to \(S_0\). For small affine parameter \(u>0\), let \(S_u \subset \mathcal N^+\) be the smooth spacelike cross‑sections with induced Riemannian metric \(h_u\), null expansion \(\theta_+(u)\), and shear tensor \(\sigma_+(u)\).  

Define the **canonical null flux functional** along \(\mathcal N^+\):  
\[
\mathcal F(u)=\int_{S_u}\!\!\Big(\sigma_+^2+\mathrm{Ric}(\ell_+,\ell_+)\Big)\,d\mu_{h_u},
\]
and let \(\mathcal A(u)=|S_u|\) denote the area of the cross-section.

**Conjecture (Null Focusing–Rigidity Inequality):**  
Under the null energy condition and smoothness assumptions above, is there a universal dimensional constant \(c_n>0\) such that  
\[
\frac{d}{du}\,\Big(\mathcal A(u)^{-\frac{1}{n-1}}\mathcal F(u)\Big)\ge c_n\,\mathcal A(u)^{-\frac{n}{n-1}}\int_{S_u}\!\theta_+^2\,d\mu_{h_u},
\]
for all sufficiently small \(u>0\)?  Moreover, if equality holds for some open interval, must the null hypersurface \(\mathcal N^+\) be totally geodesic and shear‑free with \(\mathrm{Ric}(\ell_+,\ell_+)=0\), implying a local null splitting  
\[
(\mathcal{N}^+,g|_{\mathcal{N}^+})\simeq ([0,\varepsilon)\times S_0,\,-2\,du\,dt + h_{S_0})?
\]

---

Why is it a "good" problem:  
1. **Profound insight:**  
This question articulates a potential *quantitative rigidity principle* for the null focusing mechanism.  It asks whether average null curvature and shear, rescaled by area growth, obey a canonical inequality along marginally trapped null foliations – a missing counterpart to Myers‑type and Hawking‑mass inequalities.  Proving or disproving it would clarify whether Raychaudhuri’s equation admits integral comparison forms in Lorentzian geometry.

2. **Bridge between fields:**  
It intertwines **null congruence analysis** from relativity with **comparison geometry** and the analytic structure of **spacelike mean‑curvature flows**, forging direct links among causal theory, PDE stability, and Riemannian geometric inequalities.

3. **Extreme simplicity:**  
The statement is compact and geometric: an inequality involving area, shear, Ricci focusing, and expansion, with equality implying rigidity.  Yet verifying or refuting it would demand new analytic tools extending energy estimates and monotonicity formulae to null settings.

4. **Potential to open new directions:**  
A confirmed inequality would establish the foundation of **Lorentzian comparison monotonicity theory**, governing averaged null curvatures and leading to new rigidity theorems about horizon or null‑cone structures.  Counterexamples would delineate the limits of comparison principles in degenerate (null) geometry.

Its logical coherence, explicit analytic content, and conceptual breadth make this an elegant and challenging research‑level problem.